\section{Workload and Scheduling Opportunity}
\label{sec:background}
\subsection{Sampling the target branch}
We query GitHub's merged-PR index for the named base branches of Kubernetes (\texttt{master}), TypeScript (\texttt{main}), and Next.js (\texttt{canary}). For each repository we sort by merge timestamp and PR number, retain the latest complete target window, and check its inclusive endpoints against the public index. The collector follows the PR file connection to its final page. A second query via the REST search endpoint enumerates window membership. Both interfaces depend on GitHub's index; their agreement detects acquisition mistakes but cannot establish completeness outside that index. The supplement records query boundaries, IDs, merge object IDs, file counts, and the retrieved paths. It includes the two Kubernetes PRs missed by an earlier first-parent association procedure.

Merge time is observable; queue admission is not. We use merge order to define a descriptive sequence and rescale inter-merge gaps to construct a simulated arrival process. Calling consecutive changes a window therefore describes consecutive \emph{departures}. It does not imply they were ever ready together. In particular, a rejected submission leaves no record in this corpus.

\subsection{Where packing can change a batch}
For each path we take its first two directory components, or the available prefix for a shorter path. All root files share one bucket. The FIFO column of Table~\ref{tab:trace_summary} gives the fraction of rolling eight-change windows with at least one shared bucket. The last column asks a different question: scanning a rolling sixteen-change window from its head, can a greedy packer select four changes without a shared bucket? The scan supplies a constructive opportunity under an all-ready assumption. Co-change expansion, admission constraints, and actual arrival gaps can all reduce what a scheduler obtains.

% AUTO-GENERATED BY eval/analyze_traces.py — DO NOT EDIT
\begin{table}[t]
\centering
\caption{Empirical Characteristics of the 600 Real Merged GitHub PR Traces.}
\label{tab:trace_summary}
\begin{tabular}{lccccc}
\toprule
\textbf{Repository} & \textbf{PRs} & \textbf{Uniq.} & \textbf{Gini} & \textbf{FIFO-8} & \textbf{Disj.-4} \\
 &  & \textbf{Files} &  & \textbf{Overlap} & \textbf{Feasib.} \\
\midrule
\texttt{kubernetes} & 200 & 1394 & 0.36 & 100.0\% & 90.8\% \\
\texttt{TypeScript} & 200 & 1448 & 0.30 & 100.0\% & 15.7\% \\
\texttt{next.js} & 200 & 1754 & 0.19 & 100.0\% & 91.9\% \\
\bottomrule
\end{tabular}
\end{table}


Every repository exhibits overlap in the FIFO windows, yet the opportunity to construct disjoint batches varies substantially. The TypeScript window offers fewer such batches than either of the others. A collision in this calculation means only that two paths share a directory prefix. Da Silva et al.\ independently examined changes to the same Java declarations and found that even this much finer overlap predicate admits both interfering and noninterfering edits~\cite{dasilva2024semantic}. In our queue, overlap controls scheduling priority. Treating it as an observed semantic failure would make the empirical argument circular.
